\subsection{An unattained 2-Jordan exponent: Problem 21.121(a)}
\label{cand:21.121}
\problemstatement{21.121}

A group \(\Gamma\) is \(p\)-Jordan if some constants \(J,e\) have
the following property: every finite \(H\leq\Gamma\) has a normal
abelian \(p'\)-subgroup of index at most \(J|H_{(p)}|^e\), where
\(H_{(p)}\) is a Sylow \(p\)-subgroup. The question asks whether the
infimum of the admissible exponents must itself be admissible.

\begin{proposition}
There is a countable \(2\)-Jordan group with exponent
\[
 e_0=\frac{\log24}{\log8}
\]
for which no constant works at exponent \(e_0\).
\end{proposition}

\begin{proof}
For \(k\geq1\), put \(n=2^k-1\) and index the coordinates of
\(\F_2^n\) by nonzero \(v\in\F_2^k\). Use the simplex code
\[
 C_k=\{(a\cdot v)_{v\ne0}:a\in\F_2^k\}.
\]
It has dimension \(k\). We need a small intersection estimate:
if \(L\) is spanned by \(r\) nonzero vectors with disjoint supports
and \(c=\dim(C_k\cap L)\), then
\begin{equation}\label{eq:simplex-intersection}
 2^c\leq r+1.
\end{equation}
Indeed, a basis of the intersection corresponds to \(c\) independent
linear forms on \(\F_2^k\), so its coordinate columns include all
\(2^c-1\) nonzero columns in \(\F_2^c\). Vectors in \(L\) are
constant on each of the \(r\) disjoint supports and zero elsewhere,
giving at most \(r\) different nonzero columns.

Let \(S=Q_8\rtimes C_3\), where a generator of \(C_3\) cycles
\(i,j,ij\) in \(Q_8\). Identify the center of \(Q_8^n\) with
\(\F_2^n\), and form
\[
 G_k=S^n/C_k,\qquad P_k=Q_8^n/C_k,\qquad T_k=C_3^n.
\]
The code is central and fixed by \(T_k\), so
\[
 G_k=P_k\rtimes T_k,\qquad
 |G_k|=24^n/2^k,\qquad |P_k|=8^n/2^k.
\]
Write \(Z_0\) for the image of \(Z(Q_8^n)\). Then
\(V=P_k/Z_0=\bigoplus_{i=1}^n V_i\), with each \(V_i\)
two-dimensional over \(\F_2\); the \(i\)-th \(C_3\) factor acts
irreducibly on \(V_i\) and trivially on the other coordinates.

We claim that every \(H\leq G_k\) contains a normal abelian
odd-order subgroup \(A_0\) for which, with \(u=|H_{(2)}|\),
\begin{equation}\label{eq:jordan-subgroup}
 [H:A_0]=u3^r,\qquad u\geq\frac{8^r}{r+1}
\end{equation}
for some integer \(r\geq0\).
Put \(U=H\cap P_k\), its normal Sylow \(2\)-subgroup. A Sylow
\(3\)-subgroup \(A\) complements \(U\), simply by comparing
orders. Conjugating into the Sylow subgroup \(T_k\), we may assume
\(H=U\rtimes A\) with \(A\leq T_k\). Set \(A_0=C_A(U)\);
it is normal abelian of odd order, since it centralizes \(U\)
and \(A\) is abelian. Write \(|A/A_0|=3^r\).

Let \(W\) be the image of \(U\) in \(V\). The kernel of the
action of \(A\) on \(W\) is exactly \(A_0\). To check the
nontrivial inclusion, an element \(a\) in that kernel satisfies
\(x^a=xz_x\), with \(z_x\in U\cap Z_0\), for every \(x\in U\).
It fixes \(z_x\), and \(a^3=1\), so
\(x=x^{a^3}=xz_x^3=xz_x\), whence \(z_x=1\).

The representation of \(A\) on \(W\) is completely reducible.
Over \(\F_4\), its commuting operators diagonalize, since
\(X^3-1\) has distinct roots. Nontrivial irreducibles over
\(\F_2\) therefore have dimension two and correspond to pairs
\(\{\alpha,-\alpha\}\) of nonzero homomorphisms
\(\alpha:A\to\F_3\). The occurring characters span a space of
dimension \(r\), because their common kernel is \(A_0\).
Choose \(r\) independent occurring characters and nonzero vectors
\(w_j\) in their respective isotypic components of \(W\).
Their coordinate supports are disjoint: a coordinate \(V_i\)
belongs to only one character pair.

Choose \(a_j\in A\) acting nontrivially on the character of \(w_j\)
and lift \(w_j\) to \(x_j\in U\). In each nonzero quaternion
coordinate, \(w_j\) and \(w_j^{a_j}\) are distinct nonzero vectors
of \(Q_8/Z(Q_8)\), so the commutator of their lifts is the central
involution. Thus \([x_j,x_j^{a_j}]\) is the image of the binary
support vector \(d_j\) of \(w_j\). The \(d_j\) are disjoint and
independent before quotienting by \(C_k\). If \(L\) is their
span, \eqref{eq:simplex-intersection} shows that their images span
a space of dimension at least \(r-\log_2(r+1)\) in
\([U,U]\leq U\cap Z_0\). Also \(\dim W\geq2r\).
Consequently
\[
 u=|W|\,|U\cap Z_0|
   \geq2^{2r}2^{r-c}
   \geq 8^r/(r+1).
\]
Together with \([H:A_0]=|U||A/A_0|\), this proves
\eqref{eq:jordan-subgroup}, including \(r=0\).

For \(e>e_0>1\), it follows that
\[
 \frac{[H:A_0]}{|H_{(2)}|^e}
 \leq (r+1)^{e-1}
       \left(\frac{3}{8^{e-1}}\right)^r.
\]
The right-hand side has a finite supremum over \(r\geq0\),
independent of \(k\) and \(H\), since its exponential base is
less than one.

On the other hand, every normal odd-order subgroup of \(G_k\)
centralizes \(P_k\): their commutator is in their trivial
intersection. Its image in \(T_k\) must then act trivially on
\(V\), because inner automorphisms from \(P_k\) act trivially
on \(P_k/Z_0\). The action of \(T_k\) is faithful, so that
normal subgroup is trivial. Hence the smallest permissible
index for \(G_k\) itself is \(|G_k|\), and
\[
 \frac{|G_k|}{|P_k|^{e_0}}=2^{k(e_0-1)}\longrightarrow\infty.
\]

Finally take \(\Gamma=*_{k\geq1}G_k\). Every finite subgroup
of a free product is conjugate into a factor. One proof uses
the free-product tree: a finite group fixes the center of the
finite subtree spanned by an orbit, and the action has no edge
inversions. Vertex stabilizers are trivial or conjugate to
factors. The uniform estimates therefore give every exponent
\(e>e_0\) for \(\Gamma\), while its embedded factors rule out
\(e_0\). A smaller admissible exponent would imply admissibility
of \(e_0\), so the infimum is exactly \(e_0\).
\end{proof}

The use of a free product is material; a direct sum would introduce
additional finite products not covered by the argument. No finite
generation assumption appears in the question. The construction
does not resolve part~(b). The archive contains GAP checks on all
subgroup conjugacy classes for \(k=1,2\), and direct code and
quaternion controls; the uniform argument above supplies the
unbounded quantifiers. Source: \artifact{research/21.121a-proof.md}.
